\documentclass[12pt]{article}
\usepackage{amsmath, amssymb, enumitem, fancyhdr, graphicx, libertine, nth,
  pgfplots, tikz, wasysym, xcolor}
\pgfplotsset{compat=1.11}
\usetikzlibrary{decorations.pathreplacing}
\usetikzlibrary{patterns}
\usepackage[margin=1in]{geometry}
\usepackage[libertine]{newtxmath}
\pagestyle{fancy}
\definecolor{solution-color}{HTML}{000080}
\chead{\emph{EC201 -- Practice Final -- Solutions}}

% Define new topic command
\newcommand{\topic}[1]{\medskip\begin{center}\large\textsc{#1}\normalsize
  \end{center}}

% Define question counter:
\newcounter{question}[section]
\newcommand{\question}[1]{\refstepcounter{question}
  \subsubsection*{Question~\thequestion~-- #1}}

% Subquestions are enumerate with letters
\newenvironment{subquestions}{
  \begin{enumerate}[label=(\roman*)]}{\end{enumerate}}

\begin{document}
\question{Preferences (10 Points)}
\begin{itemize}
  \item Draw a pair of indifference curves that would violate transitivity.
    Explain why they violate transitivity.

    \textcolor{solution-color}{
      Consider the diagram below. Since $A$ is on a higher indifferent curve
      than $C$, we know that $A\succ C$. $A$ and $B$ are on the same
      indifference cuve so $A\sim B$ and $B$ and $C$ are also on the same
      indifferent curve so $B\sim C$. Since $A\sim B$ and $B\sim C$,
      transitivity would say that $A\sim C$. But we just saw that $A\succ C$, a
      violation of transitivity.
    }

    \textcolor{solution-color}{
      \begin{center}
        \begin{tikzpicture}[scale=0.9]
      % Axis
          \draw [->] (0,0) node [below] {0} -- (0,0) -- (5.5,0) node [below]
          {Good 1};
          \draw [->] (0,0) node [below] {0} -- (0,0) -- (0,5.5) node [above]
          {Good 2};
      % Indifference curve
          \draw (0.3,5) to [out=280,in=175] (4.5,0.5);
          \draw (1,5) to [out=280,in=175] (3.5,0.2);
          \draw (0.3,5) node {\Large\textbullet};
          \draw (0.3,5) node [right] {$A$};
          \draw (3.5,0.19) node {\Large\textbullet};
          \draw (3.5,0.21) node [right] {$C$};
          \draw (1.53,2.1) node {\Large\textbullet};
          \draw (1.53,2.1) node [right] {$B$};
        \end{tikzpicture}
      \end{center}
    }


  \item Draw an indifference curve that violates convexity. Explain why it
    violates convexity.

    \textcolor{solution-color}{
      Here $A\sim B$. The the average bundle of $A$ and $B$ is $C$. But $C$ is
      below the indifference curve so $A\succ C$. But convexity says that
      $C\succeq A$ (the average bundle is at least as good).
    }

    \textcolor{solution-color}{
      \begin{center}
        \begin{tikzpicture}[scale=0.9]
      % Axis
          \draw [->] (0,0) node [below] {0} -- (0,0) -- (5.5,0) node [below]
          {Good 1};
          \draw [->] (0,0) node [below] {0} -- (0,0) -- (0,5.5) node
          [above] {Good 2};
      % Indifference curve
          \draw (1,5) to [out=20,in=80] (5.5,1.2);
          \draw[dashed] (1,5) -- (5.5,1.2);
          \draw (1,5) node {\Large\textbullet};
          \draw (1,5) node [left] {$A$};
          \draw (5.5,1.2) node {\Large\textbullet};
          \draw (5.5,1.2) node [right] {$B$};
          \draw (3.1,3.2) node {\Large\textbullet};
          \draw (3.1,3.1) node [above right] {$C$};
        \end{tikzpicture}
      \end{center}
    }
  \item Draw a pair of indifference curves that violate monotonicity. Explain
    why they violate monotonicity.

    \textcolor{solution-color}{
      Recall monotonicity says that if you have two bundles where one has more
      of all goods (or specificically at least as much of all goods and
      strictly more for at least one) then you should prefer that bundle. Shown
      below are two indifference curves for utility levels 1 and 2. However,
      everywhere on the indifference curve for utility level 1, you have
      more of both goods than bundles on the indifference curve for utility
      level 2, which is a violation of monotonicity.
    }
    \textcolor{solution-color}{
      \begin{center}
        \begin{tikzpicture}[scale=0.9]
        % Axis
        \draw [->] (0,0) node [below] {0} -- (0,0) -- (5.5,0) node [below]
        {Good 1};
        \draw [->] (0,0) node [below] {0} -- (0,0) -- (0,5.5) node [above]
        {Good 2};
        % Indifference curve
        \draw (0.3,5) to [out=280,in=175] (5.5,0.5);
        \draw (5.5,0.5) node [right] {$U=2$};
        \draw (1,5) to [out=280,in=175] (5.5,1.2);
        \draw (5.5,1.2) node [right] {$U=1$};
        \end{tikzpicture}
      \end{center}
    }
\end{itemize}
\question{Demand (10 Points)}
  Your utility function for goods 1 and 2 is given by:
  \[
    u\left( x_1,x_2 \right)= \log\left({x_1}\right) + x_2
  \]
\textbf{Note}: If $f\left( x \right)=\log\left( x \right)$, then
$\frac{df\left( x \right)}{dx}=\frac{1}{x}$.

\begin{subquestions}
  \item Good 1 costs \$1 and good 2 costs \$2. You have \$10 to
    spend.  How much of good 1 will you buy and how much of good 2 will you buy?

    \textcolor{solution-color}{
    \[
      MU_1 = \frac{1}{x_1}
    \]
    \[
      MU_2 = 1
    \]
    \[
      MRS = -\frac{1}{x_1}
    \]
    Setting $\left| MRS \right|=\frac{p_1}{p_2}$ gives $x_1=\frac{p_2}{p_1}$.
    Using this in the budget constraint gives:
    \[
      p_1\left( \frac{p_2}{p_1} \right)+p_2x_2 = m \;\;\;\;\Longleftrightarrow
      \;\;\;\;  x_2=\frac{m-p_2}{p_2}
    \]
    Using the prices and income we get $x_1=2$ and $x_2=4$.
  }
  \item For each good, are they normal or inferior?

    \textcolor{solution-color}{
      If income increases, it won't affect your demand for good 1. So good 1 is
      neither normal nor inferior. If income increases it will, however,
      increase your demand for good 2. Therefore good 2 is normal.
    }
\end{subquestions}

\question{Intertemporal Choice (10 Points)}
Your lifetime utility function is:
\[
u\left( c_1,c_2 \right)=c_1^{\frac{1}{2}}c_2^{\frac{1}{2}}
\]
The interest rate is 10\%. You earn income of \$100 in
period 1 and \$110 in period 2. How much will you save/borrow in period 1?

  \textcolor{solution-color}{
    The $MRS$ is:
    \[
      MRS=-\frac{\frac{1}{2}c_1^{\frac{1}{2}-1}c_2^{\frac{1}{2}}}{\frac{1}{2}c_1^{\frac{1}{2}}c_2^{\frac{1}{2}-1}}=-\frac{c_2}{
      c_1}
    \]
    The optimality condition is to set $\left| MRS \right|=\left( 1+r
    \right)$, since $-\left( 1+r \right)$ is the slope of the budget line. Therefore $c_2=\left( 1+r \right)c_1$. Using this in the
    intertemporal budget constraint:
    \begin{align*}
      c_1 + \frac{c_2}{1+r}&=m_1 + \frac{m_2}{1+r} \\
      c_1 + \frac{c_1\left( 1+r \right)}{1+r}&=100 + \frac{110}{1.1} \\
      c_1 + c_1&=100+100 \\
      2c_1 &=200 \\
      c_1&=100
    \end{align*}
    And since $c_2=\left( 1+r \right)c_1$, you consume $c_2=110$ in period 2.
    Therefore you don't borrow or save at all. You just consume your endowment.
  }

\question{Uncertainty (10 Points)}
Your total assets are worth \$30,000. Included in that \$30,000 are \$20,000
worth of valuable items in your apartment. There is a 10\% chance that you will be
burgled. The local insurance company offers to insure each dollar of stolen
items at 15\cent. If you fully insure, it will cost $\$20000\times 15\cent = \$3,000$. The remaining assets (the \$10,000) are safely locked in a
secure bank vault.

Let state 1 be the case where you are burgled and state 2 be the state where
you are not burgled. Call $c_1$ and $c_2$ your consumption in both states.
\begin{subquestions}
  \item Draw your budget contstraint for consumption across states where $c_1$
    is on the horizontal axis and $c_2$ is on the vertical axis. Assume you
    can't ``over insure'' (you can only buy a maximum of \$20,000-worth of
    insurance) and you can't negatively insure. You are, however, able to
    partially insure. Label the consumption path where you don't buy insurance
    and where you fully insure.  For each of these points, label how much you
    have to consume in each state of the world.

    \textcolor{solution-color}{
      If you don't insure, your consumption in state one is \$10,000 and
      \$30,000 in state 2. If you fully insure, your consumption is
      \$30,000-\$3,000=\$27,000 in both states.
      \begin{center}
        \begin{tikzpicture}[scale=2.2]
        % Axis
          \draw [->] (0,0) node [below] {0} -- (0,0) -- (4,0) node [below]
        {$c_1$};
        \draw [->] (0,0) node [below] {0} -- (0,0) -- (0,4) node [above]
        {$c_2$};
        \draw[thick] (1,3) -- (2.7,2.7);
        \draw[dashed] (0,0) -- (2.7,2.7);
        \draw (1,3) node {\Large{\textbullet}};
        \draw (1,3) node [above right] {No insurance};
        \draw (2.7,2.7) node {\Large{\textbullet}};
        \draw (2.7,2.7) node [above right] {Full insurance};
        \draw (0.1,0) node [above right] {$45^{\circ}$};
        \draw (1,0) node [below] {\$10,000};
        \draw[dashed] (1,0) -- (1,3);
        \draw (0,3) node [left] {\$30,000};
        \draw[dashed] (0,3) -- (1,3);
        \draw (2.7,0) node [below] {\$27,000};
        \draw[dashed] (2.7,0) -- (2.7,2.7);
        \draw (0,2.7) node [left] {\$27,000};
        \draw[dashed] (0,2.7) -- (2.7,2.7);
        \end{tikzpicture}
      \end{center}
    }
  \item What is the slope of the budget line?

    \textcolor{solution-color}{
      The slope of the budget line is
      $\frac{30,000-27,000}{10,000-27,000}=\frac{3,000}{-17,000}=-\frac{3}{17}$.
    }
  \item How much would an actuarially fair insurance policy cost?

    \textcolor{solution-color}{
      The actuarially fair insurance policy should cost $10\cent$ per dollar
      lost. To fully insure it would cost \$2,000 if it was actuarially fair.
    }

  \item What would the slope of the budget line be if the insurance was
    actuarially fair? Would it be steeper or flatter?

    \textcolor{solution-color}{
      The slope would become
      $\frac{30,000-28,000}{10,000-28,000}=\frac{2,000}{-18,000}=-\frac{1}{9}$,
      which is flatter.  Note that is the same is
      $-\frac{\gamma}{1-\gamma}=-\frac{0.1}{1-0.1}=-\frac{1}{9}$, where
      $\gamma$ is the cost per dollar insured.
    }

\end{subquestions}


\question{Equilibrium, Firm Supply and Industry Supply (12 Points)}
The demand curve for a particular market is given by:
\[
  D\left( p \right)=880-20p
\]
There are 100 firms operating in this market each with a cost function $c\left(
q \right)=\frac{q^2}{4}$.
\begin{subquestions}
  \item What is the equilibrium price and quantity?

    \textcolor{solution-color}{
      The firms will choose a quantity such that $MC\left( q \right)=p$.
      $MC\left( q \right)=c^\prime\left( q \right)=\frac{q}{2}$. Each
      individual firm's supply curve is then $S_i\left( p \right)=2p$. The
      industry supply curve is then $S\left( p \right)=100\times S_i\left( p
      \right)=100\times 2p = 200p$. We can set this supply curve equal to the
      demand curve to find the equilibrium price:
      \[
        880-20p=200p \;\;\;\;\Longleftrightarrow\;\;\;\;
        p=4
      \]
      The equilibrium quantity is then $200\times 4 = 800$.
    }
\end{subquestions}
Suddenly this good increases in popularity everyone is now willing to pay \$1
more for the good. The government decides now is a good time to introduce a
value tax of 25\% on the good.
\begin{subquestions}
  \setcounter{enumi}{1}
  \item What is the new equilibrium price and quantity as a result of the
    popularity increase and the tax? (Assume we're in the short run so that the
    number of firms does not change).

    \textcolor{solution-color}{
      To see how the increase in willingness to pay affects the demand curve we
      first need to find the inverse demand curve. This is $p\left( q
      \right)=44-\frac{1}{20}q$. The increase in the willingness to pay will
      shift this line up by one so the new intercept is 45 instead of 44. So
      the inverse demand curve is then $p\left( q
      \right)=45-\frac{1}{20}q$. Translating this back into the demand curve
      gives:
      \[
        D\left( p \right)=900-20p
      \]
      The tax will cause a wedge in between the price the buyers pay and the
      price the sellers receive. With a sales tax this is $P_B=\left( 1+\tau
      \right)P_S$. With a tax of 25\% this is $P_B=1.25 P_S$. Using this:
      \[
        900-20\times1.25 P_S = 200 P_S \;\;\;\;\Longleftrightarrow \;\;\;\; P_S=4
      \]
      With $P_S=4$, $P_B=1.25 \times 4=5$. The equilibrium quantity is
      $200\times 4 = 800$.
    }
\end{subquestions}


\question{Monopoly (13 Points)}
Suppose market demand curve for a particular good is given by $D\left( p
\right)=70-\frac{1}{3}p$. There is only one firm selling this particular good.
The cost function for this firm is $c\left( q \right)=30q+3q^2$.
\begin{subquestions}
  \item What price and quantity does the monopolist choose? What is its profits?

    \textcolor{solution-color}{
      The inverse demand curve is $p\left( q \right)=210-3q$. The revenue for
      the monopolist is $R\left( q \right)=p\left( q \right)q=\left( 210-3q
      \right)q=210q-3q^2$. The marginal revenue is then $MR\left( q
    \right)=210-6q$.
    The marginal cost is $MC\left( q \right)=c^\prime\left( q \right)=30+6q$.
    The monopolist will set $MR\left( q \right)=MC\left( q \right)$ so:
    \[
      210-6q=30+6q \;\;\;\;\Longleftrightarrow \;\;\;\;12q=180
      \;\;\;\;\Longleftrightarrow \;\;\;\; q = 15
    \]
    With this, the price is $p=210-3\times 15=165$. The profits are then:
    \[
      \pi=pq-30q-3q^2 = 165\times 15 - 30\times 15 - 3\times 15^2 =1350
    \]
    }
  \item What is the elasticity of demand at the price and quantity the
    monopolist is operating at?

    \textcolor{solution-color}{
      One way to do this is to calculate elasticity from the demand function:
      \[
        \varepsilon\left( p \right)=\frac{dq}{dp}\frac{p}{q}=-\frac{1}{3}\frac{p}{70-\frac{1}{3}p}
      \]
      At $p=165$, this is:
      \[
        \varepsilon =
        -\frac{1}{3}\frac{165}{70-\frac{165}{3}}=-3\frac{2}{3}
      \]
      Another way to find the elasticity is to know that marginal revenue can
      be expressed as $MR\left( q \right)=p\left( q \right)\left( 1+\frac{1}{
      \varepsilon\left( q \right)} \right)$. Marginal utility at $q=15$
      is $MR\left( q=15 \right)=210 - 6\times 15 = 120$. Using this in the
      formula:
      \[
        120=165\left( 1+\frac{1}{  \varepsilon} \right)
        \;\;\;\;\Longleftrightarrow \;\;\;\;  \varepsilon
        =-\frac{1}{\frac{120}{165}-1}=-3\frac{2}{3}
      \]
      The monopolist operates on the elastic part of the demand curve, as
      usual.
    }
  \item If the monopolist could perfectly price discriminate, what profits
    could it achieve?

    \textcolor{solution-color}{
      The monopolist has no fixed costs. The monopolist should charge each
      person their willingness to pay. The profits will then be the area in
      between demand and marginal cost up until demand and marginal cost
      intersect. The intersection happens at:
      \[
        210-3q= 30+6q \;\;\;\;\Longleftrightarrow \;\;\;\;
        q=20
      \]
      The price at $q=20$ is $p=210-3\times 20 = 150$.
      \begin{center}
        \begin{tikzpicture}[scale=0.9]
        % Axis
          \draw [->] (0,0) -- (10,0) node [below]
        {$q$};
        \draw [->] (0,0) -- (0,10) node [above]
        {$p$};
        \draw (0,3) -- (9,9);
        \draw (9,9) node [right] {$MC$};
        \draw (0,3) node [left] {30};
        \draw (0,9) -- (9,3);
        \draw (9,3) node [right] {$D$};
        \draw (0,9) node [left] {210};
        \draw (4.5,0) node [below] {20};
        \draw[dashed] (4.5,0) -- (4.5,6);
        \draw[dashed] (0,6) -- (4.5,6);
        \draw (0,6) node [left] {150};
        \end{tikzpicture}
      \end{center}
      The area of this triangle is $\frac{1}{2}\times 20 \times \left( 210-30
      \right)=1800$. Thus the profits the monopolist could make by perfectly
      price discriminating is 1800. This is more than just setting one price
      for all (that was 1350).
    }
\end{subquestions}

\question{Game Theory (15 Points)}
\begin{subquestions}
  \item Consider the following game:
    \setlength{\arraycolsep}{0pt}
    \[
      \begin{array}{cccccc}
        &  & \quad L \quad & & \quad R \quad  &  \\ \cline{3-6} U \quad & \vline{}
        & \quad 2,4  \quad & \vline{} & \quad 1,3 \quad & \vline{}  \\ \cline{3-6}
        M \quad & \vline{} & \quad 4,2 \quad & \vline{} & \quad 0,-1 \quad &
        \vline{} \\ \cline{3-6} D \quad & \vline{} & \quad 3,1 \quad & \vline{} &
        \quad 5,3 \quad & \vline{} \\ \cline{3-6}
      \end{array}
    \]
    \begin{itemize}
      \item Does any player have any dominated strategies?

    \textcolor{solution-color}{
      $U$ is dominated by $D$ for Player 1. No matter what player 2 does, $D$
      is always better than $U$.
    }
 
      \item Find all pure-strategy Nash equilibria of the game.

    \textcolor{solution-color}{
      There are two Nash equilibria: $(M,L)$ and $(D,R)$.
    }
    \end{itemize}
    \item Consider the following extensive form game:
          \ \\
      \begin{center}
      \begin{tikzpicture}[scale=1.5,font=\footnotesize]
      \tikzstyle{solid node}=[circle,draw,inner sep=1.5,fill=black]
      \tikzstyle{hollow node}=[circle,draw,inner sep=1.5]
      \tikzstyle{level 1}=[level distance=15mm,sibling distance=3.5cm]
      \tikzstyle{level 2}=[level distance=15mm,sibling distance=1.5cm]
      \tikzstyle{level 3}=[level distance=15mm,sibling distance=1cm]
      \node(0)[hollow node,label=above:{$P1$}]{}
      child{node[solid node,label=above left:{$P2$}]{}
      child{node[label=below:{$(5,2)$}]{} edge from parent node[left]{$A$}}
      child{node[label=below:{$(3,-1)$}]{} edge from parent node[left]{$B$}}
      child{node[label=below:{$(0,3)$}]{} edge from parent node[right]{$C$}}
      edge from parent node[left,xshift=-5]{$L$}
      }
      child{node[solid node,label=above right:{$P2$}]{}
      child{node[label=below:{$(2,2)$}]{} edge from parent node[left]{$D$}}
      child{node[label=below:{$(4,1)$}]{} edge from parent node[right]{$E$}}
      edge from parent node[right,xshift=5]{$M$}
      }
      child[grow=right]{node{$(4,0)$} edge from parent node[above]{$R$}};
      \end{tikzpicture}
      \end{center}
      Find the subgame-perfect Nash equilibirium of the game.

\textcolor{solution-color}{
The unique SPNE can be found by backward induction. Player 2 will play $C$ at the left node and $D$ at the right node. Thus, Player 1 will choose $R$. So the unique SPNE is $(R, CD)$.
}
    \end{subquestions}

\question{Oligopoly (20 Points)}
The inverse demand curve for a product is $p\left( q \right)=400-q$. Two firms
operate in this industry. Their marginal cost of production is constant at 100
and they have no fixed costs.

In this industry, the firms produce their output and then hand over all of
their output to a government distributor. The government distributor then
decides what price to charge such that all of the output will be sold. The
government distributor then gives the respective revenues back to the firms.
The government distributor does not charge a fee and has no costs.
\begin{subquestions}
  \item If the two firms compete by choosing quantities simultaneously, what
    output will each produce in equilibrium?  What will the price be? What will
    each firm's profits be?

  \textcolor{solution-color}{
      This is Cournot competition.
      In general terms, firm 1's problem is:
      \[
        \underset{q_1}{\max}\; p\left( q_1+q_2 \right)q_1 - c\left( q_1
        \right)
      \]
      The revenue term is:
      \[
        p\left( q_1+q_2 \right)q_1=\left[ 400-\left( q_1+q_2
      \right) \right]q_1= 400q_1 - q_1^2 -q_1q_2
      \]
      Therefore firm 1's problem is:
      \[
        \underset{q_1}{\max}\;  400q_1 -q_1^2 - q_1q_2 - 100q_1
      \]
      The first-order condition is then:
      \[
        400-2q_1 -q_2 -100 =0
      \]
      So $q_1\left( q_2 \right) = 150 - \frac{1}{2}q_2$. This is firm 1's best
      response function.
      For firm 2, we can do the same steps and arrive at $q_2\left( q_1
      \right) =
      150-\frac{1}{2}q_1$.
      In equilibrium, both equations will be satisfied:
      \begin{align*}
        q_1&=150-\frac{1}{2}q_2 \\
        q_2&= 150-\frac{1}{2}q_1
      \end{align*}
      Inserting the equation for $q_2$ into the equation for $q_1$:
      \begin{align*}
        q_1 &= 150-\frac{1}{2}\left( 150 - \frac{1}{2}q_1 \right) \\
        q_1 &= 150 - 75 + \frac{1}{4}q_1 \\
        \frac{3}{4}q_1 &= 75 \\
        q_1 &=100
      \end{align*}
      Using this in the best response function for firm 2 we get:
      $q_2=150-\frac{1}{2}\times100
      =100$ also. The industry as a whole then
      produces $q=2\times 100=200$. The price is then
      $p=400-200=200$.
      Each firm makes profits $\pi=\left( p-AC \right)q=\left(
      200-100\right)\times 100 = 10,000$.
  }
  \item If one firm produces first and then publicly announces how much it
    produced before the other firm decides how much to produce, how much will
    each produce? What will the price be? What will their profits be?

    \textcolor{solution-color}{
      First we need to find how the follower will react to the leader's choice
      of quantity. This is exactly how the Cournot player's react to each
      other. Therefore firm 2's reaction function is:
      \[
        q_2\left( q_1 \right)=150-\frac{1}{2}q_1
      \]
      Firm 1 then will take how firm 2 will react into account when deciding on
      a quantity. Using this expression in firm 1's profit function:
      \begin{align*}
        \pi_1\left( q_1 \right) & =  400q_1 - q_1^2 -q_1q_2-100q_1 \\
        \pi_1\left( q_1 \right) &= 300q_1 - q_1^2 -q_1\left(
        150-\frac{1}{2}q_1 \right) \\
        \pi_1\left( q_1 \right) &= 300q_1 - q_1^2 -
        150q_1-\frac{1}{2}q_1^2 \\
        \pi_1\left( q_1 \right) &= 150q_1 - \frac{1}{2}q_1^2  \\
      \end{align*}
      Taking the derivative of this and setting equal to zero:
      \[
        150- q_1=0 \;\;\;\;\Longleftrightarrow \;\;\;\; q_1 = 150
      \]
      The follower then
      will produce:
      \[
        q_2 = 150-\frac{1}{2}q_1 = 150-\frac{1}{2}\times 150=75
      \]
      The industry
      as a whole then produces $150+75=225$. The price is then $p=400-225=175$.
      The profit for the leader is $\pi_1=\left( 175-100 \right)\times 150 =
      11250$. The
      profit for the follower is $\pi_2=\left( 175-100\right)\times 75 = 5625$.
    }

  \item If both firms secretly decide together how much to produce in order to
    jointly maximize profits, how much will they produce? Assume here that if
    one firm cheats on the collusion quantity that the other firm will respond
    very aggresively, so neither firm will want to cheat on the collusion
    quantity. What will the price be? What will their profits be?

    \textcolor{solution-color}{
      In general, the problem two colluding firms face is:
      \[
        \underset{\left\{ q_1,q_2 \right\}}{\max}\; p\left( q_1+q_2
        \right)\left( q_1+q_2 \right)- c\left( q_1 \right)-c\left( q_2
      \right)
      \]
      Since the marginal costs for the two firms are constant and equal to one
      another (and there is no fixed costs), all that matters is the sum of the
      two firms' outputs. It doesn't matter if one firm
      produces all of the output, they each produce half-half, or any other
      combination. So writing the problem as choosing the industry quantity,
      $q=q_1+q_2$:
      \[
        \underset{q}{\max}\; \left( 400-q \right)q- 100q
      \]
      This is the exact same as the monopoly problem. Writing out the profit
      function more explicitly:
      \[
        \pi\left( q \right)=400q - q^2 - 100q = 300q - q^2
      \]
      The first-order condition ($\frac{d\pi\left( q \right)}{dq}=0$) is then:
      \[
        300-2q=0 \;\;\;\;\Longleftrightarrow \;\;\;\;q = 150
      \]
      The price will be $p=400-150=250$. The joint profits will be
      $\pi=300\times 150 - 150^2 = 22500$.
      If the firms split the profits 50:50, each gets 11250, which is more than
      what they would get in Cournot competition.
    }
\end{subquestions}

\end{document}
